\documentclass[10pt]{extarticle}
\usepackage[sfdefault]{FiraSans}
\usepackage[T1]{fontenc}
\usepackage[utf8]{inputenc}
\usepackage{adjustbox}
\usepackage{algorithm}
\usepackage{algorithmic}
\usepackage{amsfonts}
\usepackage{amsmath}
\usepackage{amssymb}
\usepackage[style=apa]{biblatex}
\addbibresource{references.bib}
\usepackage{booktabs}
\usepackage{breqn}
\usepackage{enumitem}
\usepackage{float}
\usepackage{geometry}
\usepackage{graphicx}
\usepackage{hyperref}
\usepackage{lipsum}
\usepackage{longtable}
\usepackage{multirow}
\usepackage{pdfpages}
\usepackage{pgfgantt}
\usepackage{setspace}
\usepackage{subcaption}
\usepackage{tabularx}
\usepackage{tikz}
\usepackage{xcolor}

\DeclareLanguageMapping{english}{english-apa}

\geometry{letterpaper, left=1in, right=1in, top=1in, bottom=1in,}

\definecolor{primary}{RGB}{0,120,215}
\definecolor{secondary}{RGB}{255,87,34}
\definecolor{background}{RGB}{245,245,245}
\definecolor{blue}{RGB}{0,62,126}

\pagecolor{background}

\hypersetup{colorlinks=true, linkcolor=primary, filecolor=secondary, urlcolor=primary, citecolor=black}

\AtBeginBibliography{\hypersetup{urlcolor=black}}

\setstretch{1.15}

\setlength{\bibhang}{0.5in}
\setlength\bibitemsep{1.5\itemsep}

\setlength{\parindent}{0pt}
\setlength{\parskip}{1em}

\nocite{*}

\begin{document}

\newcommand{\mytitlepage}[2]{
    \thispagestyle{empty}

    \begin{tikzpicture}[remember picture, overlay]
        \node [inner sep=0pt] at (current page.center) {#1}; { \node [ anchor=center, inner sep=1.25cm, rectangle, fill=blue!70!white, fill opacity=0, text opacity=1, minimum height=0.2\paperheight, minimum width=\paperwidth, text width=0.8\paperwidth, font=\fontfamily{pnc}\selectfont ] at (current page.center) {#2}; } \node [anchor=south east, outer sep=3pt] at (current page.south east) {\includegraphics[width=0.33\paperwidth]{logo.png}};
    \end{tikzpicture}

    \newpage
}

{ \mytitlepage{\includegraphics[width=\paperwidth]{background.png}}
    {
        \centering
        \fontfamily{phv}
        \vspace{-200pt} % move title up
        { \Huge
            \bfseries

            \begin{center}
                {\LARGE \textbf{Next-Generation Health Monitoring:}} \\
                {\large \textbf{A Real-Time Approach to Fitness Tracking}}
            \end{center}

            \par
        }
        \vspace{8pt}
        { \Large
            \bfseries

            Proposal

            \par
        }
        \vspace{24pt}
        {
            \begin{center}
                \begin{tabular*} {\textwidth}{@{\extracolsep{\fill}}c c c}
                    {\LARGE Jonathan Agustin} & {\LARGE Alec Anderson} & {\LARGE Brandon Smith}
                \end{tabular*}
            \end{center}
        }
    }
}

\pagenumbering{arabic}

\begin{center}
    {\LARGE \textbf{Next-Generation Health Monitoring:}} \\
    {\large \textbf{A Real-Time Approach to Fitness Tracking}}
\end{center}

\section{IoT Dataset and Application Proposal}

We propose a system that combines wearable technology, mobile application, and cloud-based analytics platform to track health and performance. The system continuously collects data, then processes it using machine learning algorithms to provide personalized health insights, including health disorders. The mobile application offers user-friendly engagement and feedback to keep the users motivated.

\section{Data Source}

We use the PMData dataset as the primary source of data for our IoT application/system. This dataset consists of health and fitness data collected from the Fitbit Versa 2 smartwatches. The original publication is located at \url{https://osf.io/preprints/osf/k2apb} (Thambawita et al., 2020). We also consider the DDXPlus dataset, which provides released patient data. This dataset has time-series data on various health conditions and symptoms. The DDXPlus dataset can be accessed at \url{https://arxiv.org/abs/2205.09148} (Fansi Tchango et al., 2022).

\section{Data Collection}

Data was collected from 16 participants over five months using Fitbit Versa 2 smartwatches. The smartwatches monitored physical activity, heart rate, and sleep patterns using sensors like the accelerometer, altimeter, and photoplethysmography. This data, combined with information such as physical activity levels, heart rate, age, height, and weight, can help users estimate the calories they burn.

\section{Observations}

The PMData dataset includes 2,440 recorded activity sessions, 20,991,392 heart rate measurements, 1,836 sleep score entries, 1,747 wellness reports, and 225 injury logs.

\section{Variables}

The objective biometrics include:

\begin{itemize}
    \item \textbf{calories.json}: Calories burned per minute.
    \item \textbf{distance.json}: Distance moved per minute (in cm).
    \item \textbf{exercise.json}: Detailed information about each activity, such as type, time, level, and performance.
    \item \textbf{heart\_rate.json}: Heart beats per minute (bpm) at a given time.
    \item \textbf{lightly\_active\_minutes.json}: Total lightly active minutes per day.
    \item \textbf{moderately\_active\_minutes.json}: Total moderately active minutes per day.
    \item \textbf{resting\_heart\_rate.json}: Resting heart rate per day.
    \item \textbf{sedentary\_minutes.json}: Total sedentary minutes per day.
    \item \textbf{sleep\_score.csv}: A score that measures the quality of sleep based on various factors, such as duration, composition, revitalization, and restlessness.
    \item \textbf{sleep.json}: A breakdown of the sleep stages (light, deep, rem, awake) per sleep.
    \item \textbf{steps.json}: Steps per minute.
    \item \textbf{time\_in\_heart\_rate\_zones.json}: Minutes spent in different heart rate zones (fat burn, cardio, peak) based on a formula of 220 minus your age.
    \item \textbf{very\_active\_minutes.json}: Total very active minutes per day.
\end{itemize}

\section{Target Users}

Our target users include casual fitness enthusiasts, professional athletes, and individuals focused on health optimization and preventive care. Our system is designed to cater to users who value a scientific and quantifiable approach to health and fitness management.

\section{Industry Fit}

Our system fits within the health and fitness industry. This industry is rapidly embracing IoT for enhanced health tracking and personalized care. Our application's ability to provide real-time insights and actionable recommendations aligns with the industry's shift towards data-driven health and wellness solutions.

\newpage

\printbibliography

\newpage

\end{document}
